\section{Complex Coordinates for Truth and Information}
\label{sec:complex-coordinates}

The Boolean support square admits a useful complex-coordinate presentation.
This presentation is conservative: it changes neither the four FDE values nor
their logical operations. It packages the two independent support coordinates
into one planar object and makes two geometric decompositions explicit.

\subsection{The raw support embedding}

For an FDE value \(v=(v^+,v^-)\in\{0,1\}^2\), define
\[
  z_{\mathrm{s}}(v)=v^+ + i v^-.
\]
The four values become
\[
  z_{\mathrm{s}}(\Nval)=0,\qquad
  z_{\mathrm{s}}(\Tval)=1,\qquad
  z_{\mathrm{s}}(\Fval)=i,\qquad
  z_{\mathrm{s}}(\Bval)=1+i.
\]
Thus the imaginary direction records negative support as an independent
coordinate; it is not identified with the numerical negative of positive
support.

Since FDE negation exchanges the two support bits, the raw coordinate obeys
\[
  z_{\mathrm{s}}(\neg v)=i\,\overline{z_{\mathrm{s}}(v)}.
\]
In the Lean implementation, the right-hand side is represented directly as
the coordinate swap \((x,y)\mapsto(y,x)\). No multiplication operation on
complex values is added to the object logic.

\begin{theorem}[Raw coordinate representation]
The map \(z_{\mathrm{s}}\) is injective. Moreover, for all FDE values \(u,v\),
\[
  \dthr(u,v)=\left|z_{\mathrm{s}}(u)-z_{\mathrm{s}}(v)\right|^2.
\]
Hence the existing combinatorial wall count is exactly squared Euclidean
distance on the four vertices of the support square.
\statusverified
\end{theorem}

\subsection{Truth/information coordinates}

The affine change of coordinates
\[
  z_{\mathrm{ti}}(v)
  =(v^+-v^-)+i(v^++v^--1)
\]
rotates and rescales the support square into a diamond. At its vertices,
\[
  z_{\mathrm{ti}}(\Tval)=1,\qquad
  z_{\mathrm{ti}}(\Fval)=-1,\qquad
  z_{\mathrm{ti}}(\Bval)=i,\qquad
  z_{\mathrm{ti}}(\Nval)=-i.
\]
The real axis therefore distinguishes truth-only from falsity-only status. The
imaginary axis distinguishes overdetermination from underdetermination: positive
imaginary direction points toward glut, negative imaginary direction toward
gap. Figure~\ref{fig:complex-coordinate-geometry} displays the two equivalent
categorical pictures.

% Two equivalent coordinate views of the four categorical FDE values.
\begin{figure}[t]
\centering
\begin{tikzpicture}[
  x=1.65cm,
  y=1.65cm,
  >=Latex,
  line cap=round,
  line join=round,
  every node/.style={font=\small}
]
  % Raw support square.
  \begin{scope}[shift={(0,0)}]
    \draw[->,thick] (-0.08,0) -- (1.28,0);
    \draw[->,thick] (0,-0.08) -- (0,1.28);
    \draw[black!45] (0,0) rectangle (1,1);

    \filldraw[fill=white,thick] (0,0) circle (2.2pt)
      node[below left] {$\Nval=0$};
    \filldraw[fill=white,thick] (1,0) circle (2.2pt)
      node[below right] {$\Tval=1$};
    \filldraw[fill=white,thick] (0,1) circle (2.2pt)
      node[above left] {$\Fval=i$};
    \filldraw[fill=white,thick] (1,1) circle (2.2pt)
      node[above right] {$\Bval=1+i$};

    \node[font=\bfseries] at (0.52,1.48) {support coordinates};
    \node at (0.52,-0.32) {$z_{\mathrm{s}}=v^+ + i v^-$};
  \end{scope}

  % Coordinate-change arrow.
  \draw[-{Latex[length=2.8mm]},very thick,blue!60!black]
    (1.45,0.55) -- (2.55,0.55);
  \node[align=center,text=blue!55!black,font=\scriptsize] at (2.00,0.28)
    {affine change};

  % Rotated truth/information diamond.
  \begin{scope}[shift={(3.85,0.55)}]
    \draw[->,thick] (-1.30,0) -- (1.34,0);
    \draw[->,thick] (0,-1.28) -- (0,1.34);
    \draw[black!45] (0,-1) -- (1,0) -- (0,1) -- (-1,0) -- cycle;

    \filldraw[fill=white,thick] (1,0) circle (2.2pt)
      node[above right] {$\Tval=1$};
    \filldraw[fill=white,thick] (-1,0) circle (2.2pt)
      node[above left] {$\Fval=-1$};
    \filldraw[fill=white,thick] (0,1) circle (2.2pt)
      node[above right] {$\Bval=i$};
    \filldraw[fill=white,thick] (0,-1) circle (2.2pt)
      node[below right] {$\Nval=-i$};

    \node[font=\bfseries] at (0,1.48) {truth/information coordinates};
    \node at (0,-1.38) {$z_{\mathrm{ti}}=(v^+-v^-)+i(v^++v^--1)$};
  \end{scope}
\end{tikzpicture}
\caption{Two equivalent embeddings of the four categorical FDE values. The
raw coordinate preserves the positive/negative support square. The transformed
coordinate separates classical truth polarity on the real axis from
glut/gap polarity on the imaginary axis. The figure represents a verified
coordinate transformation, not a new complex-valued logical connective.}
\label{fig:complex-coordinate-geometry}
\end{figure}


Negation now has the form
\[
  z_{\mathrm{ti}}(\neg v)=-\overline{z_{\mathrm{ti}}(v)}.
\]
It reverses the truth/falsity coordinate while preserving the glut/gap
coordinate. Consequently, the classical values \(1\) and \(-1\) form a
two-cycle, whereas \(i\) and \(-i\) are fixed. This is the same geometric
pattern that appears in the project's Knower fixed-point bifurcation, although
the present section proves only the general FDE coordinate identity.

\begin{theorem}[Truth/information representation]
The map \(z_{\mathrm{ti}}\) is injective, maps the four FDE values to
\(1,-1,i,-i\), and represents FDE negation by
\(z\mapsto-\overline z\).
\statusverified
\end{theorem}

\subsection{Four-cell evidence decomposition}

Consider a finite evidence profile partitioned into masses
\[
  m_T,\quad m_B,\quad m_N,\quad m_F,
\]
with total mass
\[
  m_T+m_B+m_N+m_F=m.
\]
Positive and negative support masses are
\[
  m^+=m_T+m_B,
  \qquad
  m^-=m_F+m_B.
\]
The truth/information transformation then gives the exact identity
\[
  (m^+-m^-)+i(m^++m^--m)
  =
  (m_T-m_F)+i(m_B-m_N).
\]

\begin{theorem}[Four-cell balance decomposition]
For every normalized finite four-cell profile, the real coordinate of the
transformed support masses is the strict truth/falsity balance \(m_T-m_F\),
and the imaginary coordinate is the conflict/ignorance balance \(m_B-m_N\).
\statusverified
\end{theorem}

This identity makes the role of the imaginary coordinate precise. It does not
measure contradiction alone; it measures net glut relative to gap. In
particular, two distinct four-cell distributions can agree on both balances.
The complex coordinate is therefore a projection of the richer evidence state,
not a reconstruction of it.

\subsection{A balance-sensitive glut bound}

Let \(k\) be an integer threshold and suppose both signed support masses cross
it:
\[
  k\le m_T+m_B,
  \qquad
  k\le m_F+m_B.
\]
The standard overlap argument yields \(2k\le m+m_B\). The coordinate
decomposition exposes a sharper statement.

\begin{theorem}[Exact balance-sensitive threshold-glut criterion]
A normalized finite four-cell profile is a threshold glut if and only if
\[
  2k+m_N+|m_T-m_F|\le m+m_B.
\]
The ordinary overlap bound follows by dropping the two additional nonnegative
terms.
\statusverified
\end{theorem}

Gap mass and strict truth/falsity imbalance therefore consume part of the
available threshold budget. The theorem is stronger than the coarse overlap
bound while remaining elementary and dependency-free.

\subsection{Exact complex phase regions}

Write the truth/information coordinate as (z=x+iy). At integer scale, the
two expressions
\[
  m+y+x=2m^+,
  \qquad
  m+y-x=2m^-
\]
reconstruct twice the positive and negative support masses. Comparison with
the doubled threshold therefore yields the complete phase partition:
\[
\begin{array}{rcl}
\Bval &\Longleftrightarrow& 2k\le m+y+x\ \text{and}\ 2k\le m+y-x,\\
\Tval &\Longleftrightarrow& 2k\le m+y+x\ \text{and}\ m+y-x<2k,\\
\Fval &\Longleftrightarrow& m+y+x<2k\ \text{and}\ 2k\le m+y-x,\\
\Nval &\Longleftrightarrow& m+y+x<2k\ \text{and}\ m+y-x<2k.
\end{array}
\]

\begin{theorem}[Truth/information phase classification]
Every normalized finite scaled evidence profile lies in exactly the coordinate
region corresponding to its threshold value \(\Tval,\Fval,\Bval\), or
\(\Nval\). The two region coordinates are exactly twice its signed support
masses.
\statusverified
\end{theorem}

\subsection{Multiplication, conjugation, and quarter-turn symmetry}

The presentation above deliberately withheld multiplication. The new rotation
layer equips the integer coordinate pair with Gaussian-integer multiplication
and complex conjugation. Write \(\rho(z)=iz\), so
\(\rho(x+iy)=-y+ix\).

\begin{theorem}[Coordinate operations]
For every integer coordinate \(z\),
\[
  i^2=-1,\qquad
  \mathtt{iConjugate}(z)=i\overline z,
  \qquad
  \mathtt{negConjugate}(z)=-\overline z.
\]
Four successive applications of \(\rho\) return \(z\).
\statusverified
\end{theorem}

At the categorical level, let \(\mathsf C\) denote conflation, which exchanges
\(\Bval\) and \(\Nval\) while fixing \(\Tval\) and \(\Fval\), and let
\(\rho_{\mathrm{FDE}}\) be the cycle
\[
  \Tval\longmapsto\Bval\longmapsto\Fval
  \longmapsto\Nval\longmapsto\Tval.
\]

\begin{theorem}[Conjugation and rotation]
In truth/information coordinates,
\[
  z_{\mathrm{ti}}(\mathsf C v)=\overline{z_{\mathrm{ti}}(v)},
  \qquad
  z_{\mathrm{ti}}(\rho_{\mathrm{FDE}}v)
    =i\,z_{\mathrm{ti}}(v).
\]
The categorical rotation has order four. It is not any of
\(\mathrm{id}\), \(\neg\), \(\mathsf C\), or \(\neg\circ\mathsf C\),
and for every value \(v\),
\[
  \dthr(v,\rho_{\mathrm{FDE}}v)=1.
\]
\statusverified
\end{theorem}

Negation and conflation commute and generate the identity, two reflections,
and their half turn. Adding the verified quarter turn gives the familiar
dihedral interpretation of the square. The present Lean module proves the
generating transformations and their relevant identities; it does not yet
package all eight transformations as a formal \(D_4\) group action.
\statusinterpretive

Complex notation is particularly economical here because scalar multiplication
by \(i\) composes the oriented quarter turns. The same linear action could also
be represented over \(\mathbb R^2\) by the matrix
\(\left(\begin{smallmatrix}0&-1\\1&0\end{smallmatrix}\right)\); the result
therefore establishes a canonical complex presentation rather than the
necessity of complex numbers.

\subsection{Rotation and lossless extension of four-cell evidence}

The categorical cycle lifts to profiles through
\[
  (m_T,m_B,m_N,m_F)\longmapsto(m_N,m_T,m_F,m_B).
\]

\begin{theorem}[Mass rotation and diamond containment]
For every normalized four-cell profile \(s\),
\[
  z(\rho s)=i\,z(s),
  \qquad
  |\operatorname{Re}z(s)|+|\operatorname{Im}z(s)|\le m.
\]
The formal diamond statement is given by the four equivalent linear
inequalities \(\pm\operatorname{Re}z\pm\operatorname{Im}z\le m\).
\statusverified
\end{theorem}

This is a containment result. For natural cell masses at a fixed total, parity
and integrality constrain the realized lattice points; exact surjectivity onto
the integer diamond is not asserted.

The complex projection retains two of the three independent coordinates of a
normalized four-cell profile. Define the remaining real mode by
\[
  h=(m_T+m_F)-(m_B+m_N).
\]

\begin{theorem}[Lossless extension]
For every normalized finite profile,
\[
\begin{aligned}
  4m_T&=m+h+2\operatorname{Re}z,
  &4m_F&=m+h-2\operatorname{Re}z,\\
  4m_B&=m-h+2\operatorname{Im}z,
  &4m_N&=m-h-2\operatorname{Im}z.
\end{aligned}
\]
Hence \((m,z,h)\) reconstructs the entire profile, and rotation negates the
third coordinate: \(h(\rho s)=-h(s)\).
\statusverified
\end{theorem}

The structural reading is that \(\operatorname{Re}z\) measures truth against
falsity, \(\operatorname{Im}z\) measures conflict against ignorance, and
\(h\) measures classically determined mass against glutty or gappy mass.
\statusinterpretive

\subsection{Classifier equivariance and its boundary}

\begin{theorem}[Quarter-turn equivariance away from the tie diagonal]
Let \(2k=m\). If
\(\operatorname{Im}z(s)\neq\operatorname{Re}z(s)\), then
\[
  \operatorname{val}_k(\rho s)
    =\rho_{\mathrm{FDE}}\bigl(\operatorname{val}_k(s)\bigr).
\]
\statusverified
\end{theorem}

The side condition is substantive. Threshold equality counts as support, so the
classifier's tie rule is not invariant under every quarter turn. For example,
the profile \(m_T=m_F=1\), \(m_B=m_N=0\) at \(k=1\) has \(z=0\) and value
\(\Bval\); its rotated profile again has value \(\Bval\), whereas categorical
rotation would predict \(\Fval\).

\paragraph{Formal boundary.}
Complex multiplication acts on coordinates and evidence profiles, not as a
truth function of the object language. No quantum-probabilistic interpretation
is asserted. The full dihedral group action, exact integer image of the
four-cell simplex, and model-path lift of the equivariance theorem remain open.
